\chapter{Sustitución nucleófila}\label{ch:b1:nucleophilic-substitution}

Tómese un frasco de \cip{S}-2-bromobutano, un solo enantiómero que gira
la luz polarizada. Calentado con hidróxido de sodio en un \omterm{def:b1:intermolecular-forces-solvents:solvent-classes}{disolvente
polar} aprótico, da butan-2-ol que también es un solo enantiómero --- pero
el otro, el \cip{R}: la reacción ha dado la vuelta a la molécula como el
viento a un paraguas. Dejado en cambio en agua templada, el mismo bromuro
da butan-2-ol parcialmente racémico. El mismo cambio global, un grupo OH
que sustituye a un átomo de Br, sigue dos caminos distintos. Este
capítulo establece los dos mecanismos a partir de sus \omterm{def:b1:rate-laws:order}{leyes de velocidad},
explica su estereoquímica y da las reglas que deciden entre ellos --- el
primer uso completo de las herramientas de los capítulos de cinética y
de efectos electrónicos.

\begin{recall}
Las \omterm{def:b1:rate-laws:order}{leyes de velocidad}, las \omterm{def:b1:elementary-steps:elementary-step}{etapas elementales}, la \omterm{def:b1:elementary-steps:rds}{etapa determinante de
la velocidad} y la \omterm{def:b1:elementary-steps:qssa}{aproximación del estado estacionario} están en los
\cref{ch:b1:rate-laws,ch:b1:elementary-steps}; las \omterm{def:b1:stereochemistry-in-depth:isomers}{configuraciones}, los
dibujos de Cram y la \omterm{def:b1:stereochemistry-in-depth:optical-activity}{actividad óptica}, en el
\cref{ch:b1:stereochemistry-in-depth}; los \omterm{def:b1:electronic-effects:nucleophile}{nucleófilos}, los \omterm{def:b1:electronic-effects:nucleophile}{electrófilos},
los \omterm{def:b1:electronic-effects:nucleophile}{grupos salientes} y los \omterm{def:b1:electronic-effects:intermediates}{carbocationes}, en el
\cref{ch:b1:electronic-effects}; los \omterm{def:b1:intermolecular-forces-solvents:solvent-classes}{disolventes próticos} y apróticos, en
el \cref{ch:b1:intermolecular-forces-solvents}.
\end{recall}

\section{La sustitución y sus participantes}

\begin{definition}[Sustitución nucleófila, sustrato]\label{def:b1:nucleophilic-substitution:substitution}
Una \emph{sustitución nucleófila}\index{sustitucion nucleofila@sustitución nucleófila} es la
sustitución, sobre un carbono saturado, de un \omterm{def:b1:electronic-effects:nucleophile}{grupo saliente} Y por un
\omterm{def:b1:electronic-effects:nucleophile}{nucleófilo} Nu: $\mathrm{Nu^-} + \mathrm{R{-}Y} \to \mathrm{R{-}Nu} +
\mathrm{Y^-}$. El compuesto R--Y es el \emph{sustrato}\index{sustrato};
el carbono que lleva Y es \omterm{def:b1:electronic-effects:nucleophile}{electrófilo} porque el enlace C--Y está
polarizado hacia Y.
\end{definition}

Los halogenoalcanos son los \omterm{def:b1:nucleophilic-substitution:substitution}{sustratos} típicos, con \ce{Cl-}, \ce{Br-} o
\ce{I-} como \omterm{def:b1:electronic-effects:nucleophile}{grupos salientes}. El hidróxido da alcoholes; los alcóxidos,
éteres; el cianuro, nitrilos (con un carbono más); el amoniaco y las
aminas, aminas; el yoduro se intercambia con los demás haluros. El agua y
los alcoholes, usados como disolventes, también pueden actuar como
\omterm{def:b1:electronic-effects:nucleophile}{nucleófilos}: la sustitución es entonces una \omterm{def:b1:nucleophilic-substitution:sn1}{solvólisis}.

\section{El mecanismo SN2}

Cuando el bromometano reacciona con hidróxido en agua, duplicar
cualquiera de las dos concentraciones duplica la velocidad:
$v = k[\ce{CH3Br}][\ce{OH-}]$, una reacción de orden uno respecto de cada
reactivo.

\begin{definition}[Mecanismo SN2, inversión de Walden]\label{def:b1:nucleophilic-substitution:sn2}
El \emph{mecanismo SN2}\index{mecanismo SN2} (\omterm{def:b1:nucleophilic-substitution:substitution}{sustitución nucleófila}
bimolecular) es una sola \omterm{def:b1:elementary-steps:elementary-step}{etapa elemental} en la que el \omterm{def:b1:electronic-effects:nucleophile}{nucleófilo} ataca el
carbono por el lado opuesto al \omterm{def:b1:electronic-effects:nucleophile}{grupo saliente} mientras este se va. Los
otros tres sustituyentes del carbono se dan la vuelta, como un paraguas
que el viento vuelve del revés: la \emph{inversión de
Walden}\index{inversion de Walden@inversión de Walden}.
\end{definition}

\begin{omfigure}
\setchemfig{atom sep=2.1em}
\schemestart
\chemfig{H@{nu}O^{-}}
\arrow{0}[,0.35]
\chemfig{@{c}C(-[:150]H_3C)(<[:210]H)(<:[:250]C_2H_5)-[@{cb}]@{br}Br}
\arrow{->}
\chemfig{[:0]HO-C(-[:30]CH_3)(<[:-30]H)(<:[:-70]C_2H_5)}
\+
\chemfig{Br^{-}}
\schemestop
\chemmove{\draw[->, thick, shorten <=2pt, shorten >=3pt](nu) .. controls +(15:5mm) and +(175:5mm) .. (c);
          \draw[->, thick, shorten <=2pt](cb) .. controls +(90:6mm) and +(90:6mm) .. (br);}

\medskip
\begin{tikzpicture}[font=\small]
  \node[draw, dashed, rounded corners, inner sep=6pt] at (0,0)
    {$\left[\vcenter{\hbox{\ce{HO}$\cdots$\chemfig{C(-[2]CH_3)(-[:-30]H)(-[:210]C_2H_5)}$\cdots$\ce{Br}}}\right]^{-}$};
  \node[below] at (0,-1.1) {estado de transición: cinco grupos alrededor del C, tres en un plano};
\end{tikzpicture}

{\small La reacción SN2 del hidróxido con el \cip{S}-2-bromobutano. El
\omterm{def:b1:electronic-effects:nucleophile}{nucleófilo} ataca la cara opuesta al bromo (\omterm{def:b1:electronic-effects:curly-arrow}{flechas curvas}); en el \omterm{def:b1:elementary-steps:energy-profile}{estado
de transición}, el enlace O--C se está formando mientras se rompe el
enlace C--Br; el producto es \cip{R}-butan-2-ol: los otros tres grupos se
han dado la vuelta hacia el otro lado.}
\end{omfigure}

\begin{proposition}[Ley de velocidad de la SN2]\label{prop:b1:nucleophilic-substitution:sn2-law}
Para una reacción SN2, $v = k[\mathrm{RY}][\mathrm{Nu}]$.
\end{proposition}

\begin{proof}
El mecanismo es una sola \omterm{def:b1:elementary-steps:elementary-step}{etapa elemental} bimolecular; por la ley de
acción de masas de las \omterm{def:b1:elementary-steps:elementary-step}{etapas elementales} (\cref{ch:b1:elementary-steps}),
su velocidad es proporcional al producto de las concentraciones de los
dos participantes.
\end{proof}

\begin{proposition}[Estereoquímica de la SN2]\label{prop:b1:nucleophilic-substitution:sn2-inversion}
Una reacción SN2 sobre un carbono estereogénico da un solo
estereoisómero, con la disposición invertida de los tres grupos que no
cambian: una reacción de este tipo es estereoespecífica.
\end{proposition}

\begin{proof}
El \omterm{def:b1:electronic-effects:nucleophile}{nucleófilo} solo puede llegar al lado libre del carbono, opuesto al
\omterm{def:b1:electronic-effects:nucleophile}{grupo saliente}; los otros tres grupos pasan por un plano en el \omterm{def:b1:elementary-steps:energy-profile}{estado de
transición} y terminan al otro lado. El enlace nuevo apunta adonde no
apuntaba el antiguo: cada molécula de \omterm{def:b1:nucleophilic-substitution:substitution}{sustrato} da la disposición
especular de sus grupos. El descriptor suele cambiar de \cip{S} a
\cip{R} o a la inversa, pero no siempre, ya que depende de las
prioridades del \omterm{def:b1:electronic-effects:nucleophile}{nucleófilo} y del \omterm{def:b1:electronic-effects:nucleophile}{grupo saliente}.
\end{proof}

\begin{definition}[Reacción estereoespecífica]\label{def:b1:nucleophilic-substitution:stereospecific}
Una reacción es una \emph{reacción estereoespecífica}\index{reaccion estereoespecifica@reacción estereoespecífica}
cuando \omterm{def:b1:nucleophilic-substitution:substitution}{sustratos} \omterm{def:b1:stereochemistry-in-depth:isomers}{estereoisómeros} dan productos \omterm{def:b1:stereochemistry-in-depth:isomers}{estereoisómeros} distintos,
cada \omterm{def:b1:nucleophilic-substitution:substitution}{sustrato} el suyo.
\end{definition}

\section{El mecanismo SN1}

El 2-bromo-2-metilpropano reacciona con el agua a una velocidad que no
depende de la concentración del \omterm{def:b1:electronic-effects:nucleophile}{nucleófilo}: $v = k[(\ce{CH3})_3\ce{CBr}]$.
Una etapa en la que no interviene el \omterm{def:b1:electronic-effects:nucleophile}{nucleófilo} debe decidir la
velocidad.

\begin{definition}[Mecanismo SN1, racemización, solvólisis]\label{def:b1:nucleophilic-substitution:sn1}
El \emph{mecanismo SN1}\index{mecanismo SN1} (\omterm{def:b1:nucleophilic-substitution:substitution}{sustitución nucleófila}
unimolecular) tiene dos etapas: la \omterm{def:b1:electronic-effects:cleavage}{heterólisis} lenta del enlace C--Y,
que da un \omterm{def:b1:electronic-effects:intermediates}{carbocatión}, y después la adición rápida del \omterm{def:b1:electronic-effects:nucleophile}{nucleófilo} al
\omterm{def:b1:electronic-effects:intermediates}{carbocatión}. Una reacción que convierte un solo enantiómero en una
mezcla de los dos es una \emph{racemización}\index{racemizacion@racemización}; una sustitución
en la que el \omterm{def:b1:electronic-effects:nucleophile}{nucleófilo} es el disolvente es una
\emph{solvólisis}\index{solvolisis@solvólisis}.
\end{definition}

\begin{omfigure}
\setchemfig{atom sep=2.1em}
\schemestart
\chemfig{C(-[2]CH_3)(-[4]H_3C)(-[6]CH_3)-[@{cb2}]@{br2}Br}
\arrow{->[lenta]}
\chemfig{C^{+}(-[2]CH_3)(-[:210]H_3C)(-[:-30]CH_3)}
\+
\chemfig{Br^{-}}
\schemestop

\smallskip
\schemestart
\chemfig{C^{+}(-[2]CH_3)(-[:210]H_3C)(-[:-30]CH_3)}
\arrow{->[rápida][\ce{H2O}]}
\chemfig{C(-[2]CH_3)(-[4]H_3C)(-[6]CH_3)-OH_2^{+}}
\arrow{->[\ce{-H+}]}
\chemfig{C(-[2]CH_3)(-[4]H_3C)(-[6]CH_3)-OH}
\schemestop
\chemmove{\draw[->, thick, shorten <=2pt](cb2) .. controls +(90:6mm) and +(90:6mm) .. (br2);}

\medskip
\begin{tikzpicture}[font=\small]
  \fill[omProp!15, draw=omProp] (0,0.2) ellipse (0.2 and 0.85);
  \fill[omProp!15, draw=omProp] (0,-0.2) ellipse (0.2 and 0.85);
  \draw[thick] (0,0) -- (-1.1,-0.2) node[left] {\ce{CH3}};
  \draw[thick] (0,0) -- (0.95,-0.45) node[right] {\ce{C2H5}};
  \draw[thick] (0,0) -- (0.4,0.5) node[above right] {H};
  \node[fill=white, inner sep=1pt] at (0,0) {\ce{C+}};
  \draw[->, thick, omDef] (-0.9,1.6) node[left] {\ce{H2O}} -- (-0.15,0.95);
  \draw[->, thick, omDef] (0.9,-1.6) node[right] {\ce{H2O}} -- (0.15,-0.95);
  \node[align=center] at (4.6,0) {ataque por cualquier cara:\\ las dos configuraciones,\\ racémico si nada más\\ distingue las caras};
\end{tikzpicture}

{\small \omterm{def:b1:nucleophilic-substitution:sn1}{Solvólisis} SN1 del 2-bromo-2-metilpropano en agua: ionización
lenta al \omterm{def:b1:electronic-effects:intermediates}{carbocatión} terciario y, después, captura por el agua y pérdida
de un protón. Abajo, un \omterm{def:b1:electronic-effects:intermediates}{carbocatión} secundario \omterm{def:b1:stereochemistry-in-depth:stereocentre}{quiral} (procedente del
2-bromobutano) es plano, y el agua se adiciona por cualquiera de sus
caras con la misma probabilidad.}
\end{omfigure}

\begin{proposition}[Ley de velocidad de la SN1]\label{prop:b1:nucleophilic-substitution:sn1-law}
Para el \omterm{def:b1:nucleophilic-substitution:sn1}{mecanismo SN1}, si el \omterm{def:b1:electronic-effects:intermediates}{carbocatión} se captura mucho más deprisa de
lo que vuelve al \omterm{def:b1:nucleophilic-substitution:substitution}{sustrato}, $v = k_1[\mathrm{RY}]$, independiente del
\omterm{def:b1:electronic-effects:nucleophile}{nucleófilo}.
\end{proposition}

\begin{proof}
Etapas: $\mathrm{RY} \rightleftharpoons \mathrm{R^+} + \mathrm{Y^-}$ ($k_1$,
$k_{-1}$), $\mathrm{R^+} + \mathrm{Nu} \to \mathrm{RNu}$ ($k_2$). El
\omterm{def:b1:electronic-effects:intermediates}{carbocatión} es un \omterm{def:b1:electronic-effects:intermediates}{intermedio reactivo}: por la \omterm{def:b1:elementary-steps:qssa}{aproximación del estado
estacionario} (\cref{ch:b1:elementary-steps}), $k_1[\mathrm{RY}] =
k_{-1}[\mathrm{R^+}][\mathrm{Y^-}] + k_2[\mathrm{R^+}][\mathrm{Nu}]$, de
modo que
$v = k_2[\mathrm{R^+}][\mathrm{Nu}] = \dfrac{k_1k_2[\mathrm{RY}][\mathrm{Nu}]}
{k_{-1}[\mathrm{Y^-}] + k_2[\mathrm{Nu}]}$. Cuando $k_2[\mathrm{Nu}] \gg
k_{-1}[\mathrm{Y^-}]$, se reduce a $k_1[\mathrm{RY}]$: la primera etapa es
la determinante.
\end{proof}

\begin{proposition}[Estereoquímica de la SN1]\label{prop:b1:nucleophilic-substitution:sn1-stereo}
Una reacción SN1 sobre un carbono estereogénico da las dos
\omterm{def:b1:stereochemistry-in-depth:isomers}{configuraciones}: la reacción no es estereoespecífica y conduce a la
\omterm{def:b1:nucleophilic-substitution:sn1}{racemización} cuando nada distingue las dos caras del \omterm{def:b1:electronic-effects:intermediates}{carbocatión}.
\end{proposition}

\begin{proof}
El \omterm{def:b1:electronic-effects:intermediates}{carbocatión} es plano (\cref{ch:b1:electronic-effects}): sus dos caras
son imágenes especulares. Si el \omterm{def:b1:electronic-effects:nucleophile}{grupo saliente} se ha alejado antes de la
captura, el \omterm{def:b1:electronic-effects:nucleophile}{nucleófilo} se adiciona a cualquiera de las caras a la misma
velocidad y da los dos \omterm{def:b1:stereochemistry-in-depth:enantiomers}{enantiómeros} en cantidades iguales. En la
práctica, el ion saliente sigue tapando una cara durante un tiempo, y a
menudo se observa un ligero exceso de inversión.
\end{proof}

\begin{omfigure}
\begin{tikzpicture}
\begin{axis}[
  width=0.8\linewidth, height=5.0cm, xmin=0, xmax=10, ymin=-30, ymax=110,
  axis x line=bottom, axis y line=left, xlabel={coordenada de reacción},
  ylabel={energía}, xtick=\empty, ytick=\empty,
  label style={font=\small}, legend style={font=\scriptsize, draw=none, at={(0.98,0.98)}, anchor=north east}]
\addplot[omDef, very thick] table[x=x, y=E] {figdata/out/bachelor-1-energy-profiles-sn2.dat};
\addplot[omProp, very thick, dashed] table[x=x, y=E] {figdata/out/bachelor-1-energy-profiles-sn1.dat};
\legend{SN2: una etapa, SN1: dos etapas}
\node[font=\scriptsize, anchor=north] at (axis cs:5,68) {carbocatión};
\node[font=\scriptsize, anchor=north] at (axis cs:0.6,-3) {RY};
\node[font=\scriptsize, anchor=south] at (axis cs:9.6,-18) {RNu};
\end{axis}
\end{tikzpicture}

{\small \omterm{def:b1:elementary-steps:energy-profile}{Perfiles de energía} esquemáticos. SN2: un solo \omterm{def:b1:elementary-steps:energy-profile}{estado de
transición}, con cinco grupos alrededor del carbono. SN1: una primera
barrera, más alta (la ionización, la etapa determinante), el \omterm{def:b1:electronic-effects:intermediates}{carbocatión}
en un pozo poco profundo y, después, una pequeña barrera hasta su
captura.}
\end{omfigure}

\section{Qué decide el mecanismo}

\begin{proposition}[El sustrato]\label{prop:b1:nucleophilic-substitution:substrate}
Los \omterm{def:b1:nucleophilic-substitution:substitution}{sustratos} metílicos y primarios reaccionan por SN2; los terciarios,
por SN1; los secundarios, por cualquiera de los dos, según los demás
factores. Los \omterm{def:b1:nucleophilic-substitution:substitution}{sustratos} alílicos y bencílicos reaccionan deprisa por
ambos.
\end{proposition}

\begin{proof}
La SN2 necesita espacio detrás del carbono: cada grupo alquilo en lugar
de un hidrógeno abarrota el \omterm{def:b1:elementary-steps:energy-profile}{estado de transición}, que lleva cinco grupos,
y frena la reacción; tres grupos la bloquean. La SN1 necesita un
\omterm{def:b1:electronic-effects:intermediates}{carbocatión} estable: el orden de estabilidad
(\cref{prop:b1:electronic-effects:carbocation-order}) hace la ionización
rápida en los terciarios, posible en los secundarios y demasiado lenta
para los \omterm{def:b1:electronic-effects:intermediates}{carbocationes} primarios y metílicos. Los \omterm{def:b1:electronic-effects:intermediates}{carbocationes} alílicos
y bencílicos se estabilizan por resonancia, y el mismo sistema $\pi$
estabiliza el \omterm{def:b1:elementary-steps:energy-profile}{estado de transición} de la SN2.
\end{proof}

\begin{proposition}[El grupo saliente]\label{prop:b1:nucleophilic-substitution:leaving-group}
Los dos mecanismos son tanto más rápidos cuanto mejor es el \omterm{def:b1:electronic-effects:nucleophile}{grupo
saliente}, es decir, cuanto más débil es la base en la que se convierte:
$\ce{I-} > \ce{Br-} > \ce{Cl-} \gg \ce{F-}$; \ce{OH-} y \ce{RO-} no se
van.
\end{proposition}

\begin{proof}
En ambos, el enlace C--Y se rompe en la etapa determinante o antes de
ella, y el \omterm{def:b1:elementary-steps:energy-profile}{estado de transición} lleva parte de la carga negativa sobre Y.
Un grupo que retiene bien una carga negativa --- la base conjugada de un
\omterm{def:b1:predominant-reaction:strong-acid}{ácido fuerte} --- rebaja ese \omterm{def:b1:elementary-steps:energy-profile}{estado de transición}. \ce{HI}, \ce{HBr} y
\ce{HCl} son \omterm{def:b1:predominant-reaction:strong-acid}{ácidos fuertes}; \ce{HF}, uno débil ($\mathrm{p}K_a$ 3.2); el
agua y los alcoholes, muy débiles: \ce{OH-} y \ce{RO-} son \omterm{def:b1:predominant-reaction:strong-acid}{bases fuertes} y
malos \omterm{def:b1:electronic-effects:nucleophile}{grupos salientes}.
\end{proof}

\begin{proposition}[El disolvente]\label{prop:b1:nucleophilic-substitution:solvent}
Los \omterm{def:b1:intermolecular-forces-solvents:solvent-classes}{disolventes polares} próticos (agua, alcoholes) favorecen la SN1; los
\omterm{def:b1:intermolecular-forces-solvents:solvent-classes}{disolventes polares} apróticos (propanona, dimetilsulfóxido,
dimetilformamida) favorecen la SN2 con \omterm{def:b1:electronic-effects:nucleophile}{nucleófilos} aniónicos.
\end{proposition}

\begin{proof}
La SN1 crea dos iones a partir de un \omterm{def:b1:nucleophilic-substitution:substitution}{sustrato} neutro: un \omterm{def:b1:intermolecular-forces-solvents:solvent-classes}{disolvente
polar} prótico estabiliza ambos, el catión mediante sus \omterm{def:b1:lewis-resonance-vsepr:lewis}{pares no
enlazantes} y el anión mediante \omterm{def:b1:intermolecular-forces-solvents:hbond}{enlaces de hidrógeno}, y rebaja la barrera
de ionización. En la SN2, la carga de un \omterm{def:b1:electronic-effects:nucleophile}{nucleófilo} aniónico se reparte
en el \omterm{def:b1:elementary-steps:energy-profile}{estado de transición}; un \omterm{def:b1:intermolecular-forces-solvents:solvent-classes}{disolvente prótico}, que une fuertemente
el anión pequeño mediante \omterm{def:b1:intermolecular-forces-solvents:hbond}{enlaces de hidrógeno}, la frena, mientras que un
\omterm{def:b1:intermolecular-forces-solvents:solvent-classes}{disolvente aprótico} deja el anión libre y reactivo
(\cref{ch:b1:intermolecular-forces-solvents}).
\end{proof}

\begin{method}[Elegir el mecanismo]\label{met:b1:nucleophilic-substitution:choose-mechanism}
\begin{enumerate}
  \item \omterm{def:b1:nucleophilic-substitution:substitution}{Sustrato}: metílico o primario $\to$ SN2; terciario $\to$ SN1;
        secundario $\to$ se sigue.
  \item \omterm{def:b1:electronic-effects:nucleophile}{Nucleófilo}: fuerte y aniónico ($\ce{OH-}$, $\ce{RO-}$, $\ce{CN-}$,
        $\ce{I-}$) $\to$ SN2; débil y neutro (agua, alcohol) $\to$ SN1.
  \item Disolvente: aprótico $\to$ SN2; prótico $\to$ SN1.
  \item Se predice la estereoquímica (inversión o \omterm{def:b1:nucleophilic-substitution:sn1}{racemización}) y se
        comprueba que una \omterm{def:b1:predominant-reaction:strong-acid}{base fuerte} no provoque más bien una
        eliminación
        (\cref{ch:b1:elimination}).
\end{enumerate}
\end{method}

\begin{omfigure}
\begin{tabular}{lll}
\hline
 & SN2 & SN1 \\
\hline
etapas & una & dos, a través de un \omterm{def:b1:electronic-effects:intermediates}{carbocatión} \\
\omterm{def:b1:rate-laws:order}{ley de velocidad} & $k[\mathrm{RY}][\mathrm{Nu}]$ & $k[\mathrm{RY}]$ \\
\omterm{def:b1:nucleophilic-substitution:substitution}{sustrato} & metilo $>$ primario $>$ secundario & terciario $>$ secundario \\
\omterm{def:b1:electronic-effects:nucleophile}{nucleófilo} & fuerte, aniónico & débil, neutro (a menudo el disolvente) \\
disolvente & polar aprótico & polar prótico \\
estereoquímica & inversión (estereoespecífica) & \omterm{def:b1:nucleophilic-substitution:sn1}{racemización} (en general) \\
\hline
\end{tabular}

\smallskip
{\small Los dos mecanismos, uno junto a otro.}
\end{omfigure}

\section{Ejercicios}

\begin{exercise}[$\star$]\label{exo:b1:nucleophilic-substitution:1}
Escríbanse los productos de: el bromoetano con hidróxido de sodio; el
yodometano con etóxido de sodio; el 1-cloropropano con cianuro de
potasio; el bromometano con amoniaco (primera etapa). Identifíquense el
\omterm{def:b1:nucleophilic-substitution:substitution}{sustrato}, el \omterm{def:b1:electronic-effects:nucleophile}{nucleófilo} y el \omterm{def:b1:electronic-effects:nucleophile}{grupo saliente}.
\end{exercise}

\begin{exercise}[$\star$]\label{exo:b1:nucleophilic-substitution:2}
Para una reacción $\mathrm{RY} + \mathrm{Nu}$ con
$v = k[\mathrm{RY}][\mathrm{Nu}]$, ¿qué le ocurre a la \omterm{def:b1:rate-laws:initial-rate}{velocidad inicial}
si se triplica $[\mathrm{Nu}]$? ¿Y si se dividen por dos ambas
concentraciones? Las mismas preguntas si $v = k[\mathrm{RY}]$.
\end{exercise}

\begin{exercise}[$\star$]\label{exo:b1:nucleophilic-substitution:3}
Predígase el mecanismo (SN1 o SN2): el 1-bromobutano con yoduro de sodio
en propanona; el 2-bromo-2-metilpropano en agua; el bromometano con
hidróxido; el 2-bromobutano en etanol sin base añadida.
\end{exercise}

\begin{exercise}[$\star$]\label{exo:b1:nucleophilic-substitution:4}
Dibújese en \omterm{def:b1:stereochemistry-in-depth:representations}{representación de Cram} el producto de la reacción SN2 del
cianuro con el \cip{R}-2-bromobutano y dese su descriptor.
\end{exercise}

\begin{exercise}[$\star\star$]\label{exo:b1:nucleophilic-substitution:5}
Escríbase el \omterm{def:b1:nucleophilic-substitution:sn1}{mecanismo SN1} completo de la hidrólisis del
2-bromo-2-metilpropano, con \omterm{def:b1:electronic-effects:curly-arrow}{flechas curvas} e incluida la pérdida del
protón. ¿Por qué la \omterm{def:b1:rate-laws:order}{ley de velocidad} es de primer orden aunque
intervengan tres especies?
\end{exercise}

\begin{exercise}[$\star\star$]\label{exo:b1:nucleophilic-substitution:6}
Explíquese por qué el 1-bromo-2,2-dimetilpropano, un bromuro primario,
reacciona muy despacio por SN2 y no reacciona por SN1.
\end{exercise}

\begin{exercise}[$\star\star$]\label{exo:b1:nucleophilic-substitution:7}
El \cip{S}-2-yodooctano pierde su \omterm{def:b1:stereochemistry-in-depth:optical-activity}{actividad óptica} cuando se mantiene en
propanona con yoduro de sodio, aunque no se forma ningún compuesto nuevo.
Explíquese, y dígase por qué la velocidad de pérdida de actividad es el
doble de la velocidad de sustitución.
\end{exercise}

\begin{exercise}[$\star\star$]\label{exo:b1:nucleophilic-substitution:8}
Ordénense por velocidad de SN2 con yoduro: el 2-bromo-2-metilpropano, el
bromometano, el 2-bromopropano y el 1-bromopropano. Ordénense los mismos
compuestos por velocidad de \omterm{def:b1:nucleophilic-substitution:sn1}{solvólisis} SN1.
\end{exercise}

\begin{exercise}[$\star\star$]\label{exo:b1:nucleophilic-substitution:9}
Un bromuro secundario ópticamente puro sufre una \omterm{def:b1:nucleophilic-substitution:sn1}{solvólisis}; el producto
gira la luz al \qty{12}{\percent} del valor esperado para una inversión
pura. Dense las proporciones de los dos \omterm{def:b1:stereochemistry-in-depth:enantiomers}{enantiómeros} del producto e
interprétense.
\end{exercise}

\begin{exercise}[$\star\star\star$]\label{exo:b1:nucleophilic-substitution:10}
Dedúzcase la \omterm{def:b1:rate-laws:order}{ley de velocidad} general de la SN1 con la \omterm{def:b1:elementary-steps:qssa}{aproximación del
estado estacionario}, y demuéstrese que añadir al principio el ion del
\omterm{def:b1:electronic-effects:nucleophile}{grupo saliente}, \ce{Y-}, frena la reacción (el \omterm{def:b1:precipitation:common-ion}{efecto del ion común} en
cinética). ¿En qué condición es despreciable este efecto?
\end{exercise}

\begin{exercise}[$\star\star\star$]\label{exo:b1:nucleophilic-substitution:11}
Dos \omterm{def:b1:rate-laws:half-life}{tiempos de semirreacción} para la reacción del bromometano con
hidróxido: con $[\ce{OH-}]_0 = [\ce{CH3Br}]_0 = C_0$, demuéstrese que
$1/[\ce{CH3Br}] = 1/C_0 +
kt$ y dese $t_{1/2}$. Con $[\ce{OH-}]_0 = 50\,C_0$, demuéstrese que el
decrecimiento es exponencial y dese $t_{1/2}$. Datos del ejercicio:
$k =
\qty{2.0e-4}{L/(mol.s)}$, $C_0 = \qty{0.010}{mol/L}$.
\end{exercise}

\begin{exercise}[$\star\star\star$]\label{exo:b1:nucleophilic-substitution:12}
Explíquese por qué un alcohol, \ce{ROH}, no reacciona con el bromuro de
sodio, pero sí con el ácido bromhídrico concentrado, para dar \ce{RBr}.
Escríbase el mecanismo para un alcohol terciario y para uno primario.
\end{exercise}

\section{Problema: la solvólisis del cloruro de terc-butilo}

\begin{problem}[{Problema de fin de semana --- el orden y la constante de
velocidad a partir de medidas de conductividad, el tiempo de
semirreacción, el mecanismo y su estereoquímica, y la energía de
activación de la solvólisis}]
\label{pb:b1:nucleophilic-substitution:1}
El 2-cloro-2-metilpropano (cloruro de terc-butilo) se disuelve, a baja
concentración, en una mezcla de agua y propanona. Su \omterm{def:b1:nucleophilic-substitution:sn1}{solvólisis}
\ce{(CH3)3CCl + H2O -> (CH3)3COH + H+ + Cl-} produce iones, y se registra
la conductividad $\kappa$ de la disolución, que es proporcional a la
cantidad de iones formada. Al cabo de mucho tiempo alcanza
$\kappa_\infty = \qty{2.000}{mS/cm}$ a las dos temperaturas. Medidas
(\unit{mS/cm}):

\begin{center}
\begin{tabular}{lcccccccc}
\hline
$t$ (min) & 0 & 5 & 10 & 15 & 20 & 30 & 45 & 60 \\
\hline
$\kappa$, \qty{25.0}{\celsius} & 0 & 0.172 & 0.329 & 0.473 & 0.605 & 0.835 & 1.110 & 1.321 \\
$\kappa$, \qty{35.0}{\celsius} & 0 & 0.507 & 0.885 & 1.168 & 1.379 & 1.654 & 1.856 & 1.940 \\
\hline
\end{tabular}
\end{center}
% data generated in figdata/bachelor-1/solvolysis.py (story data)

\medskip
\noindent\textbf{Parte I --- El orden.}
\begin{enumerate}
  \item ¿Por qué la conductividad es proporcional a la cantidad de
        \omterm{def:b1:nucleophilic-substitution:substitution}{sustrato} que ha reaccionado?
  \item Exprésese $[\mathrm{RCl}]/[\mathrm{RCl}]_0$ en función de $\kappa$ y
        $\kappa_\infty$.
  \item Escríbanse la ley integrada de una reacción de primer orden y la
        magnitud que hay que representar frente a $t$ para comprobarla.
  \item Calcúlese $\ln(\kappa_\infty - \kappa)$ a \qty{25.0}{\celsius} en
        cada instante.
  \item ¿Es la representación una recta? Conclúyase sobre el orden.
  \item El agua está en gran exceso. ¿Qué oculta eso en el orden?
\end{enumerate}

\medskip
\noindent\textbf{Parte II --- \omterm{def:b1:rate-laws:order}{Constante de velocidad} y \omterm{def:b1:rate-laws:half-life}{tiempo de semirreacción}.}
\begin{enumerate}[resume]
  \item Hállese la \omterm{def:b1:rate-laws:order}{constante de velocidad} a \qty{25.0}{\celsius}, en
        \unit{s^{-1}}.
  \item Calcúlese el \omterm{def:b1:rate-laws:half-life}{tiempo de semirreacción}.
  \item ¿En qué instante se ha consumido el \qty{90}{\percent} del
        \omterm{def:b1:nucleophilic-substitution:substitution}{sustrato}?
  \item Hállese la \omterm{def:b1:rate-laws:order}{constante de velocidad} a \qty{35.0}{\celsius}.
  \item ¿En qué factor ha aumentado la velocidad con diez grados?
  \item Compruébese un valor de la serie de \qty{35}{\celsius} frente a la
        ley.
\end{enumerate}

\medskip
\noindent\textbf{Parte III --- El mecanismo.}
\begin{enumerate}[resume]
  \item ¿Qué mecanismo sugieren el \omterm{def:b1:nucleophilic-substitution:substitution}{sustrato} y el disolvente?
  \item Escríbase con \omterm{def:b1:electronic-effects:curly-arrow}{flechas curvas}.
  \item Demuéstrese que concuerda con el orden hallado.
  \item Añadir hidróxido de sodio a \qty{0.01}{mol/L} apenas cambia la
        \omterm{def:b1:rate-laws:rate}{velocidad de desaparición} del cloruro. ¿Por qué?
  \item El mismo experimento con el cloruro terciario \omterm{def:b1:stereochemistry-in-depth:stereocentre}{quiral}
        3-cloro-3-metilhexano, ópticamente puro, da un alcohol de
        \omterm{def:b1:stereochemistry-in-depth:optical-activity}{actividad óptica} casi nula. Explíquese.
  \item Dibújese un \omterm{def:b1:elementary-steps:energy-profile}{perfil de energía} de la reacción y márquese la \omterm{def:b1:elementary-steps:rds}{etapa
        determinante de la velocidad}.
  \item ¿Por qué el 1-clorobutano apenas reaccionaría en las mismas
        condiciones?
\end{enumerate}

\medskip
\noindent\textbf{Parte IV --- La \omterm{def:b1:rate-laws:activation-energy}{energía de activación}.}
\begin{enumerate}[resume]
  \item Escríbase la \omterm{thm:b1:rate-laws:arrhenius}{ley de Arrhenius}.
  \item Exprésese $E_a$ a partir de las dos \omterm{def:b1:rate-laws:order}{constantes de velocidad}.
  \item Calcúlese $E_a$ ($R = \qty{8.314}{J/(K.mol)}$).
  \item ¿Qué mide esta energía, en el mecanismo?
  \item Predígase la \omterm{def:b1:rate-laws:order}{constante de velocidad} a \qty{45.0}{\celsius}.
  \item Dese la \omterm{def:b1:rate-laws:activation-energy}{energía de activación} de la \omterm{def:b1:nucleophilic-substitution:sn1}{solvólisis}, con dos cifras
        significativas.
\end{enumerate}
\end{problem}
